%    Include other referenced packages here.
\usepackage{pgfplots}
\usepackage{mathtools}
\usepackage{amssymb,amsmath,amsthm}
%\usepackage[inline]{enumitem}
\usepackage{relsize}
\usepackage{mathtools}
\usepackage{multicol}
\usepackage{booktabs}
\usepackage{calc}
\usepackage{esdiff}
\usepackage{systeme}
\usepackage{bm}
\usepackage[most]{tcolorbox}

%%%%%%%%%%%%%%%%%%%%%%%%%%%%%%%%%%%%%%%%%%%%%%%%%%%%%%%%%%%%%%%%%%%%%%%%
% definitions
% theorem, proposition, corollary, lemma
% \newtheorem{theorem}{Theorem}[section]
% \newtheorem{corollary}{Corollary}[theorem]
% \newtheorem{lemma}[theorem]{Lemma}
% \newtheorem{proposition}{Proposition}[section]
% \newtheorem{definition}{Definition}[section]


\newcommand{\mydef}[2]{
    \begin{tcolorbox}[colback=blue!5, colframe=blue!75!black, title=#1]
        #2
    \end{tcolorbox}
}


% numerical sets
\newcommand{\real}{{\mathbb R}}
\newcommand{\realm}{{{\mathbb R}^m}}
\newcommand{\rational}{{\mathbb Q}}
\newcommand{\integer}{{\mathbb Z}}
\newcommand{\naturals}{{\mathbb N}}
\newcommand{\complex}{{\mathbb C}}
\newcommand{\F}{{\mathbb F}}
\newcommand{\R}{{\mathbb R}}
\newcommand{\C}{{\mathbb C}}
\newcommand{\Rn}{{\mathbb R}^n}
\newcommand{\Rm}{{\mathbb R}^m}
\newcommand{\cn}{{\mathbb C}^n}
\newcommand{\rnn}{{{\mathbb R}^{n\times n}}}
\newcommand{\rnm}{{{\mathbb R}^{n\times m}}}
\newcommand{\rmn}{{{\mathbb R}^{m\times n}}}
\newcommand{\cnn}{{{\mathbb C}^{n\times n}}}

% abs and norm, ceil and floor, inner product
\DeclarePairedDelimiter\abs{\lvert}{\rvert}%
\DeclarePairedDelimiter\norm{\lVert}{\rVert}%
\DeclarePairedDelimiter\ceil{\lceil}{\rceil}
\DeclarePairedDelimiter\floor{\lfloor}{\rfloor}
\DeclarePairedDelimiter\inner{\langle}{\rangle}
%\providecommand{\norm}[1]{\lVert#1\rVert}
% Swap the definition of \abs* and \norm*, so that \abs
% and \norm resizes the size of the brackets, and the 
% starred version does not.
\makeatletter
\let\oldabs\abs
\def\abs{\@ifstar{\oldabs}{\oldabs*}}
%
\let\oldnorm\norm
\def\norm{\@ifstar{\oldnorm}{\oldnorm*}}
\makeatother

% common functions
\newcommand{\rtr}{\real\to\real}
\newcommand{\rttr}{\real^2\to\real}

% common matrices
\newcommand{\MmnF}{M_{mn}(\F)}
\newcommand{\MnF}{M_{n}(\F)}

% diagonal matrix
\DeclareMathOperator{\diag}{diag}

%domain 
\DeclareMathOperator{\dom}{dom}

% image of a linear operator
\DeclareMathOperator{\Ima}{Im}

% range of a linear operator
\DeclareMathOperator{\range}{range}

% Rank of a linear operator
\DeclareMathOperator{\rank}{rank}

% projection of a linear operator
\DeclareMathOperator{\proj}{proj}

% Null and Column space of a matrix
\DeclareMathOperator{\rspace}{\mathcal{R}}
\DeclareMathOperator{\nspace}{\mathcal{N}}
\DeclareMathOperator{\cspace}{\mathcal{C}}

% Trace of a linear operator
\DeclareMathOperator{\tr}{tr}

% Rank of a linear operator
\DeclareMathOperator{\vspan}{span}

% Sign of a function
\DeclareMathOperator{\sign}{sgn}

% differential d
\newcommand{\dif}{\mathop{}\!\mathrm{d}}
\newcommand{\od}[2]{\frac{\dif #1}{\dif #2}}
\newcommand{\dod}[2]{\displaystyle\frac{\dif #1}{\dif #2}}

% complement
\newcommand{\scomp}[1]{{#1^\complement}}

% interior & closure
\newcommand{\interior}[1]{{\kern0pt#1}^{\mathrm{o}}}
\newcommand{\closure}[1]{{\overline{#1}}}

% evaluated integral 
\newcommand\eval[3]{\left[#1\right]_{#2}^{#3}}

% matrices p q
\newcommand{\rmat}[2]{{{\mathbb R}^{{#1}\times {#2}}}}

% vectors
\newcommand{\vect}[1]{\mathbf{#1}}
\newcommand{\va}{\mathbf{a}}
\newcommand{\vb}{\mathbf{b}}
\newcommand{\vc}{\mathbf{c}}
\newcommand{\vd}{\mathbf{d}}
\newcommand{\ve}{\mathbf{e}}
\newcommand{\vf}{\mathbf{f}}
\newcommand{\vp}{\mathbf{p}}
\newcommand{\vq}{\mathbf{q}}
\newcommand{\vr}{\mathbf{r}}
\newcommand{\vs}{\mathbf{s}}
\newcommand{\vx}{\mathbf{x}}
\newcommand{\vy}{\mathbf{y}}
\newcommand{\vz}{\mathbf{z}}
\newcommand{\vu}{\mathbf{u}}
\newcommand{\vw}{\mathbf{w}}
\newcommand{\vv}{\mathbf{v}}
\newcommand{\vT}{\mathbf{T}}
\newcommand{\vX}{\mathbf{X}}
\newcommand{\vY}{\mathbf{Y}}

\newcommand{\vzero}{\mathbf{0}}
\newcommand{\mzero}{\mathbf{O}}

\newcommand{\vtn}[4]{
\vect{#1}=\left[
\begin{array}{c}
#2 \\
#3 \\
#4
\end{array}
\right]}

\newcommand{\vtu}[3]{
\left[
\begin{array}{c}
#1 \\
#2 \\
#3
\end{array}
\right]}

\newcommand{\vtd}[4]{
\left[
\begin{array}{c}
#1 \\
#2 \\
#3 \\
#4
\end{array}
\right]}

\newcommand{\vdn}[3]{
\vect{#1}=\left[
\begin{array}{c}
#2 \\
#3 
\end{array}
\right]}

\newcommand{\vdu}[2]{
\left[
\begin{array}{c}
#1 \\
#2 
\end{array}
\right]}

\newcommand{\stlb}[3]{
\left\{\begin{matrix}
#1 \\ #2 \\ #3
\end{matrix}\right.}

\newcommand{\stlt}[3]{
\left\{\begin{matrix}
#1 \\ #2 \\ #3
\end{matrix}\right.}

% horbars in matrices
\newcommand{\brows}[1]{%
  \begin{bmatrix}
  \begin{array}{@{\protect\rotvert\;}c@{\;\protect\rotvert}}
  #1
  \end{array}
  \end{bmatrix}
}
\newcommand{\rotvert}{\rotatebox[origin=c]{90}{$\vert$}}
\newcommand{\rowsvdots}{\multicolumn{1}{@{}c@{}}{\vdots}}

